\subsection{Broader literature assessment and contribution scope}\label{sec:novelty}
The starting contribution is the explicit, certified failure of the published original-rate triplet implication. Sections~\ref{sec:orderedconductances}--\ref{sec:degenerateextension} add a surviving all-dimension exclusion theorem and a finite four-state asymptotic classification; their precise novelty is assessed separately from the established tools used in their proofs. Section~\ref{sec:mechanism} also supplies a controlled family and a proved onset classification for a specified eigenvector geometry. This is a problem-specific structural result; its individual publication priority has not been established. The supporting results organize and apply established mathematics. The following comparison was checked on 11 October 2026; it separates mathematical validity from publication priority.

\paragraph{The matrix criterion and reset algebra have direct precedents.}
\href{https://arxiv.org/abs/2007.11433v3}{Baake and Sumner, \emph{On equal-input and monotone Markov matrices}, v3, 25 September 2022}, Fact 3.1 and Proposition 3.11, give the cumulative-coordinate characterization of stochastic monotonicity for matrices and generators. In their row-generator convention, let $T_{ij}=\mathbf1_{i\geq j}$. Its coordinates are $(f_1,Bf)$, so
\[
T^{-1}LT=\begin{pmatrix}0&*\\0&K\end{pmatrix}.
\]
The upper row outside the diagonal consists of outgoing-rate tail sums and is automatically nonnegative. Their transformed-generator criterion is therefore exactly our condition that $K$ is Metzler. Our Perron and covariance argument is an application of that order structure, not a new monotonicity criterion.

Their Section 2.1 also treats equal-input generators $C-cI$, with identical rows in $C$ and $C^2=cC$. Our reset generator is the transpose of this construction, with $C=\mathbf1\pi_b^\top$ and $c=1$. Its scalar action on nonstationary modes is elementary equal-input algebra. The useful calculation here is its compatibility with the particular triplet inequalities, $F(R_0)=G(R_0)=0$; the general reset construction is not claimed as new.

\paragraph{The divergence corollary also follows from established thermomajorization.}
\href{https://arxiv.org/html/2410.06686v2}{Vu and Hayakawa, \emph{Thermomajorization Mpemba Effect}, v2, 6 January 2025}, give the Lorenz-curve and all-convex-divergence characterizations of thermomajorization. The connection to our corollary is explicit. In cooling, both $p_h/\pi_b$ and $p_c/\pi_b$ are increasing, so both curves use the order $n,n-1,\ldots,1$. At their common horizontal knots,
\[
x_k=\sum_{i=k}^n\pi_{bi},\qquad
y_{h,k}=\sum_{i=k}^np_{hi}\geq\sum_{i=k}^np_{ci}=y_{c,k},
\]
by stochastic ordering. Linear interpolation orders the entire curves: $p_h$ thermomajorizes $p_c$. Heating uses the order $1,\ldots,n$ and lower-tail inequalities, so $p_c$ thermomajorizes $p_h$. This yields the corollary from the known convex-divergence characterization. The supporting-line proof above is an alternative elementary derivation; neither derivation establishes priority for a new general theorem.

This absence argument does not assert a thermomajorization reversal for our counterexample. Vu and Hayakawa's Result 4 distinguishes that stronger notion from the scalar slow-overlap criterion. The counterexample's stated conclusions remain eventual fixed-norm and KL ordering, with separate finite-time KL and TV brackets.

\paragraph{What is claimed, and what is retained only as support.}
\begin{itemize}
\item \textbf{Central result:} the four-state witness satisfying all original-rate triplet failures, its exact certification and the consequent refutation of $(\mathcal C_4)$. The checked sources do not supply this witness or settle that precise implication. This bounded finding is not an exhaustive priority guarantee.
\item \textbf{Established tools:} reversible symmetrization, slow-mode overlaps, Perron theory, stochastic monotonicity, likelihood-ratio evolution, thermomajorization, and equal-input reset algebra. The report claims no methodological novelty for these tools or for rational certification itself.
\item \textbf{Applications proved here:} the additional rate inequalities, the all-$\kappa\leq1$ subfamily, the nonincreasing gap-law subfamily, and the all-time convex-divergence corollary. Their role is to explain sufficient conditions that survive the failed conjecture. They are not promoted as independently novel theorems.
\item \textbf{Problem-specific deductions:} the identity $F=hG+dH$, the feasible-cone formulation, and reset compatibility simplify construction and verification. They are proved directly; their individual publication priority remains unestablished. The open-neighborhood result is a continuity consequence, not a new perturbation principle.
\end{itemize}

The review covers the exact source conjecture, OHI's coupling result and SI, Fill--Kahn's comparison and density-ratio argument, Baake--Sumner's matrix and generator criteria, and Vu--Hayakawa's thermomajorization formulation. Searches combined Mpemba, triplet, monotonicity, convex divergence, likelihood ratio, nonincreasing rates and equal-input generators. Relevant full-text sections were inspected, rather than inferred from titles. No identical all-lower-branch or gap-law statement was located in this bounded review; that absence is insufficient evidence of novelty. Unindexed work and exhaustive citation tracing remain outside the check. The counterexample remains a contribution independently of supporting priority questions. The added contrast-cone and classification results must be evaluated on their exact hypotheses and conclusions; the bounded search is not a guarantee that these consequences are unprecedented.

\paragraph{Additional sources for the classification and realization arguments.}
The explicit slow-mode realization is already proved in \href{https://arxiv.org/abs/1711.05829v3}{Klich et al., v3, 18 December 2018}, Appendix VIII, particularly Eq.~(67). Our rank-one formula is that construction in different normalization, and the rank-$r$ version uses the same projector idea. The Mpemba index and parity criteria are likewise prior results from that work. \href{https://arxiv.org/abs/2502.01758v3}{Teza et al., v3, 16 April 2025}, Sections 2.1 and 5.2.3, review those results. Both primary full texts were inspected on 11 October 2026. The university-hosted journal PDF failed to download; the arXiv review version was inspected instead, and differences in the final journal text have not been excluded.

The nonvanishing-Wronskian characterization of extended complete Chebyshev systems is classical; Gasull, Geyer and Ma\~nosas (2020), Lemma 2.3, states the criterion used here. The Gibbs Wronskian formula and the convex-curve test are explicit applications. The weighted Helmert basis, Metzler semigroup positivity and Perron theory are also established. What is supplied here is the rate-to-contrast identity, the resulting asymmetric cooling/heating implications of the triplet premise, and the explicit finite thermal test. Searches combining Mpemba with Helmert, conductance ordering, strong-effect absence, Chebyshev systems and stationary-point counts did not locate matching statements in the inspected primary sources. This is a bounded attribution assessment, not an exhaustive priority claim.

\paragraph{An explicit prior strong-temperature count.}
A further search located \href{https://doi.org/10.5281/zenodo.21298061}{Parker
Emmerson, \emph{Organizational Mechanics of Effective Cardinality Transitions}},
Zenodo v1, with a stated publication date of 10 July 2026 and repository creation
date of 17 July 2026. The full text was inspected on 11 October 2026.
Theorem 7.3 applies the classical sign-variation bound to the slow-overlap
numerator; Corollary 7.5 subtracts the forced bath zero. Thus the scalar
nonbath strong-temperature bound $n-2$ is explicitly present in prior work
and is not a novelty claim here. Section 11 gives a four-state sharpness example
using a preparation landscape different from the stationary one. Warning 7.13
distinguishes simultaneous mode cancellation from the scalar result.
Our rank-$r$ bound $n-r-1$ addresses that simultaneous-cancellation problem;
its priority remains subject to further review. The bound on \emph{stationary
points} and the derivative-curve classifier concern different objects from
zeros of the overlap numerator. No statement about the manuscript's peer-review
status is needed for this attribution.

The coexistence of strong cooling and heating in four states is also prior
work: \href{https://arxiv.org/abs/2308.04557v1}{Bera, Walker and Vucelja,
v1, 8 August 2023}, Section IV and Fig.~4, give a four-state ring example and
discuss nonoverlap of the strong effects in three states. Those passages
were checked in the primary full text on 11 October 2026. Our rational
coexistence example is a sharpness and verification construction, not the
first observation of that phenomenon. The general derivative test also
covers ordinary effects that need not cancel the slow mode.

\paragraph{The sharp full-triplet conclusions.}
The additional results in Sections~\ref{sec:fulltripletbound} and
\ref{sec:tripletendpoint} concern the precise original-rate premise:
a sharp $25\%$ amplitude-advantage supremum and an exact endpoint covariance
criterion for ordinary cooling. Their proofs use a restriction on the first
three components of the slow eigenfunction, established through a reversible
Schur complement and three nonnegative-rate boundary faces. A separate
Perron-coupling proof supports the same restriction. These are proposed
problem-specific contributions, not new Schur-complement or Perron principles.
Searches combining Mpemba with amplitude bounds, triplet constraints,
four-state systems and energy-span thresholds located no matching theorem
in the inspected primary sources. Expert assessment of priority remains
necessary. The spacing-independent exclusion also explains why the March
paper's finite-temperature sampling box lies entirely within a region where
the original inference is valid.
